    \chapter{Introdução}
    \label{chap:introduction}

    \begin{introduction}
    Este capítulo contextualiza e motiva o trabalho, ao apresentar o computador como fonte de dados para monitorização contínua do estado do utilizador, apresenta o foco principal e exemplos de casos de uso para a \textit{framework} proposta e define os objetivos esperados para este trabalho.
    \end{introduction}

    \section{Contexto e Motivação}

    O uso prolongado do computador em contexto profissional ou estudantil cria oportunidades únicas para a monitorização contínua do estado do utilizador, visto que permite não só o uso de sensores tradicionais (como sensores autónomos ou os sensores embutidos num dispositivo \textit{wearable}), mas também aproveitar os próprios periféricos de entrada do computador (como o rato, o teclado ou a \textit{webcam}).
    Estes dispositivos, embora não sejam tipicamente classificados como sensores, possibilitam a extração de dados biométricos e comportamentais capazes de refletir o estado do utilizador, permitindo, por exemplo, inferir o estado emocional com base na dinâmica de digitação no teclado~\cite{keyboard_speed_study,keyboard_correction_study}.

    Este tipo de monitorização é particularmente útil em contextos não presenciais como o ensino remoto ou o teletrabalho, onde a distância física limita a perceção do estado do utilizador~\cite{online_classes_sharma}, bem como em situações presenciais onde a observação direta poderá não ser suficiente, como a deteção de stresse em empregados.
    
    Contudo, ao contrário dos estudos laboratoriais, onde os dados apenas interessam no final da recolha, num cenário real o objetivo é responder às mudanças do estado de cada indivíduo à medida que estas ocorrem.
    Por isso, apesar de ser possível processar um fluxo contínuo utilizando \textit{batching} - ao agrupar os dados em lotes com janelas temporais fixas~\cite{Akidau2018} - esta abordagem introduz um \textit{delay} entre o evento e a sua visualização~\cite{ibm_stream_processing}.
    Em contrapartida, uma abordagem baseada em \textit{streaming} processa cada evento da fonte de dados de forma independente e imediata, reduzindo drasticamente o \textit{delay} entre a ocorrência do evento e a sua visualização e permitindo a análise destes dados em tempo real, surgindo assim como alternativa natural para cenários de monitorização contínua.

    \section{Cenário Representativo}

    De modo a orientar a escolha dos objetivos, bem como o desenvolvimento da prova de conceito e a sua demonstração, estabelece-se também um cenário representativo concreto: a recolha de dados de três sensores (teclado, câmara e \textit{smartwatch} WearOS) de um utilizador e o cálculo das métricas relevantes em tempo real.
    
    Contudo, este cenário assume-se estritamente como um exemplo representativo e não como uma restrição estrita de âmbito. 
    Por esta razão, algumas das decisões foram tomadas considerando a aplicabilidade da prova de conceito a outros cenários de monitorização contínua, como o suporte à monitorização de vários utilizadores ou à adição de novos sensores.

    \section{Objetivos}\label{sec:objectives}

    O objetivo deste trabalho é mostrar que é possível conjugar uma solução técnica baseada na monitorização simultânea de utilizadores durante o uso normal do computador, recorrendo apenas ao uso de periféricos normalmente presentes - interfaces do próprio computador e/ou dispositivos \textit{wearable off-the-shelf}.

    Para concretizar esta meta global, definem-se os seguintes objetivos específicos:
    \begin{itemize}
        \item \textbf{Verificar a adequação de uma abordagem baseada em \textit{streaming}:} demonstrar, com base nos requisitos típicos do cenário acima descrito - nomeadamente, baixa latência, desacoplamento entre fontes de dado e mecanismos de ingestão e suporte a fontes de dados heterogéneas - que uma arquitetura baseada em \textit{streaming} constitui uma solução válida para este tipo de problema.
        \item \textbf{Demonstrar o fluxo dos dados e o funcionamento em tempo real:} rastrear uma mensagem ao longo do seu percurso na solução e verificar de forma empírica a produção de métricas relevantes para o teclado, câmara e \textit{smartwatch} WearOS, observando a sua reatividade com base nas ações do utilizador.
        \item \textbf{Suportar a extensão futura:} Apesar do cenário de referência se restringir a um único utilizador e com um conjunto fixo de sensores, é importante garantir que a arquitetura do sistema suporta a monitorização de múltiplos utilizadores e a adição futura de novos sensores.
    \end{itemize}

    
    \section{Estrutura do Documento}
    Este documento está organizado em cinco capítulos. O Capítulo~\ref{chap:introduction} descreve o contexto e motivação, apresente o cenário principal e casos de uso e define os objetivos base do trabalho.
    O Capítulo~\ref{chap:estadodaarte} apresenta possíveis cenários de utilização, com base na literatura existente, descreve as fontes de dados utilizadas nestes cenários e apresenta as alternativas existentes e as suas lacunas.
    O Capítulo~\ref{chap:arquitetura} descreve de que forma os objetivos influenciaram a arquitetura, apresenta a arquitetura lógica e descreve as tecnologias escolhidas para a concretizar.
    O Capítulo~\ref{chap:resultados} descreve os resultados obtidos, as escolhas tomadas durante a implementação e discute os pontos fortes e fracos da solução apresentada como prova de conceito.
    Finalmente, o Capítulo~\ref{chap:conclusions} sumariza o trabalho realizado e sugere direções para trabalhos futuros.


